\def\stat{l2-14}





\def\tit{МОДЕЛЬ ИМПАКТ-ФАКТОРА В~СИСТЕМЕ ПРАВОВОГО РЕГУЛИРОВАНИЯ 


ДЕЯТЕЛЬНОСТИ НАУЧНЫХ ОРГАНИЗАЦИЙ}








\def\titkol{Модель импакт-фактора в системе правового регулирования 


деятельности %научных 


организаций} 





\def\autkol{Г.\,В.~Лукьянов}





\def\aut{Г.\,В.~Лукьянов$^1$}





\titel{\tit}{\aut}{\autkol}{\titkol}





%{\renewcommand{\thefootnote}{\fnsymbol{footnote}}\footnotetext[1]


%{Работа выполнена при поддержке РФФИ, грант 10-01-00480.}}





\renewcommand{\thefootnote}{\arabic{footnote}}


\footnotetext[1]{Институт проблем информатики 


Российской  академии наук, gena-mslu@mail.ru}      


    


    \Abst{Рассмотрен новый показатель результативности 


деятель\-ности научных организаций, а именно импакт-фактор публикаций 


работников научной организации в системе <<Web of Science>> (ИФПР). Доказано, 


что этот показатель в его новом толковании в лучшем случае сводится к 


среднему значению традиционного импакт-фактора нескольких журналов, а 


в худшем~--- к им\-пакт-фактору одного журнала.}


    


    \KW{импакт-фактор; рецензируемый журнал; оценка деятельности 


научных организаций; количество публикаций; цитирование; ссылка; 


публикационная активность}





       \vskip 12pt plus 9pt minus 6pt








      \thispagestyle{headings}





      


            \label{st\stat}


            


            \vspace*{-9pt}





    


\section{Введение}


    


    Попытки создать систему оценки деятельности научных организаций 


предпринимались не раз и в нашей стране, и за рубежом, тем более что 


показатели ее результативности зародились и сформировались в среде 


научного сообщества и фактически уже много лет применяются для тех или 


иных оценок. В~частности, количество публикаций в ряде случае выступает 


одним из важнейших показателей деятельности ученых. Например, для 


предъявления диссертации к защите в России действует правило, в 


соответствии с которым от соискателя ученой степени, наряду с прочим, 


требуется некоторое минимальное количество публикаций.


    


    В этом смысле не стала неожиданностью и Типовая методика оценки 


результативности деятельности научных организаций, выполняющих 


на\-уч\-но-ис\-сле\-до\-ва\-тель\-ские, опыт\-но-кон\-струк\-тор\-ские и технологические работы 


гражданского назначения (ТМО), введенная в действие приказом министра 


образования и науки Российской Федерации от 14.10.2009 №\,406~[1].


    


    Конечно, здесь много можно говорить о точности и корректности 


формулировок, щепетильности выбора конкретных показателей и даже о 


соотношении их значимости для обобщенной оценки результативности 


деятельности научных организаций. Однако, если закрыть глаза на все эти 


<<мелочи>>, то в ТМО все же появилась до сих пор неизвестная 


<<изюминка>>, а именно показатель с не совсем ясными предназначением и 


методикой расчета. Речь в данном случае идет об импакт-факторе 


публикаций работников научной организации в <<Web of 


Science>>\footnote{Web of Science~--- информационный ресурс в 


Интернете, предоставляемый компанией Thomson Scientific, который 


включает в том числе и сведения об импакт-факторах примерно 8600 


журналов, издаваемых во всем мире.}.


    


    Дальнейшее рассмотрение никак не связано с выбором системы 


цитирования, а только со смыслом самого показателя. Действительно, 


почему выбрана именно <<Web of Science>>, а не одна из значимых систем 


цитирования Китая как второй экономики мира, остается вопросом 


политическим. Значение же показателя~--- это содержательный вопрос, т.\,е.\ 


поддающийся научному изучению.





\section{Методика расчета импакт-фактора}





    Как известно, импакт-фактор в его оригинальной разработке 


применяется для оценки значимости журналов\footnote[2]{Именно журналов, и 


никак иначе этот термин не воспринимался вплоть до появления ТМО.} в 


произвольном году в сопоставлении с другими из некоторого выбранного 


круга журналов. Исходно этот показатель был введен Институтом научной 


информации США\footnote[3]{Institute for Scientific Information~--- коммерческая организация, 


образованная в 1960~г.\ Юджином Гарфилдом на основе созданной ранее компании Eugene Garfield 


Associates Inc. В~1992~г.\ Институт приобретен корпорацией Thomson и ныне называется Thomson 


Scientific.  Занимается составлением библиографических баз данных научных публикаций, их 


индексированием и определением индекса цитируемости, импакт-фактора и других 


статистических показателей научных работ.} в 1960-х~гг., он 


вычисляется следующим образом:


    \begin{equation}


    I_n^i=\fr{C_{n-(1,2)}^i}{B^i_{n-(1,2)}}\,,


    \label{1-2-14}


    \end{equation}


    где $I_n^i$~--- импакт-фактор $i$-го журнала в $n$-м году;


$ C_{n-(1,2)}^i$~--- чис\-ло цитирований в течение $n$-го года в 


выбранном круге журналов статей, опубликованных в $i$-м журнале в $(n-1)$-м и $(n-2)$-м годах; 


$B^i_{n-(1,2)}$~--- число статей, опубликованных в $i$-м журнале 


в $(n-1)$-м и $(n-2)$-м годах.


    


    Например, для вычисления импакт-фактора журнала 


<<Scientometrics>>\footnote[4]{Журнал в области наукометрии, 


выпускается непрерывно с 1978~г. В~настоящее время 


издается совместно Akad\'{e}miai Kiad\'{o} и Springer Science\;+\;Business 


Media (ISSN 0138-9130).} в 2009~г.\ подсчитывают число цитирований в 


2009~г.\ во всех журналах системы <<Web of Science>> статей, 


опубликованных в журнале <<Scientometrics>> в 2007 и 2008~гг.\, и 


полученное значение делят на количество статей, опубликованных в журнале 


<<Scientometrics>> в 2007 и 2008~гг. В~результате получается число 


2,17~\cite{2-2-14},  которое означает, что в среднем на каждую статью 


журнала <<Scientometrics>> в 2007~г.\ приходится 2,17 ссылки во всех 


журналах, учитываемых в системе <<Web of Science>>.


    


    Эта оценка отдаленно напоминает среднюю температуру по больнице, 


так как в журнале с высоким импакт-фактором, например, может оказаться 


статья, на которую никто не ссылался в искомом году. Тем не менее, автор 


этой статьи находится в равном <<положении>> с авторами других статей с 


точки зре\-ния им-\linebreak пакт-фак\-то\-ра. В~то же время в другом году число ссылок на 


эту кон\-крет\-ную\linebreak статью может значительно превысить этот показатель для 


других ста\-тей с прежним результатом с точки зрения импакт-фактора.


    


    Другая проблема импакт-фактора состоит в том, что традиционно 


журналы в одной научной области имеют более высокий импакт-фактор по 


сравнению с журналами из других областей. Эти два обстоятельства могут 


сформировать крайне ложное представление о ценности журнала. Для 


сравнения, им\-пакт-фак\-тор российского журнала <<Атомная энергия>> 


(Atomic Energy) в том же 2009~г.\ составил 0,062~\cite{3-2-14}, т.\,е.\ в 


35~раз меньше, чем у журнала <<Scientometrics>> в 2009~г. Трудно поверить 


по этим конкретным данным, что значимость журнала <<Атомная энергия>> 


в~35~раз ниже, чем <<Scientometrics>>.


    


    У импакт-фактора много других недостатков, главный из которых 


состоит в том, что он не отражает качество исследования, изложенного в той 


или иной статье. Наиболее понятна его коммерческая составляющая, которая 


заключается в том, что инвестор научного исследования в той или иной 


научной области может захотеть сравнить результаты исследователей для 


оценки перспектив своих инвестиций. Для этого, в частности, и 


используются объективные численные показатели, такие как импакт-фактор. 


Поэтому на подобные измерения и сущест-\linebreak вует спрос.





\section{Содержание импакт-фактора автора}


    


    Можно предположить, что разработчики ТМО, ясно понимая недостатки 


традиционного содержания и предназначения импакт-фактора, предложили 


некоторое его существенное улучшение. Здесь, правда, вызывает удивление 


тот факт, что вместо описания ясной и понятной аналитической 


функциональной зависимости математическими символами алгоритм 


вычисления ИФПР приводится в ТМО в виде текста. Такой подход 


отбрасывает нас в эпоху еще до Франсуа Виета\footnote{Франсу\'{a} 


Ви\'{e}т (фр.\ Fran\c{c}ois Vi\`{e}te, seigneur de la Bigoti\`{e}re; 


1540\,--\,13~февраля 1603)~--- французский математик, основоположник 


символической алгебры. По образованию и основной профессии~--- юрист. 


Наиболее известна теорема Виета, посвященная соотношению между 


корнями и коэффициентами уравнения второй степени.}, когда все 


выражения и уравнения описывали именно таким образом.


    


    Судя по лаконичным разъяснениям, приведенным в ТМО, ключевым 


элементом вычисления ИФПР является публикационная активность научного 


сотрудника в каждом журнале системы <<Web of Science>> за выбранный 


год (ПАС), которая подсчитывается следующим образом:


    \begin{equation}


    A^i_{j-n}=P^i_{j-n} I_n^i\,,


    \label{e2-2-14}


    \end{equation}


где $ A^i_{j-n}$~--- ПАС $j$-го сотрудника в $i$-м журнале за $n$-й 


год;  $ P^i_{j-n}$~--- число публикаций $j$-го сотрудника 


в $i$-м журнале за $n$-й год; $I_n^i$~--- импакт-фактор 


$i$-го журнала в $n$-м году.


    


    Тогда искомый показатель ИФПР, указанный в ТМО, вычисляется 


следующим образом:


    \begin{equation}


    \mathrm{ИФПР} =\fr{\sum\limits_{n=t-6}^{t-1}\left[ 


\sum\limits_{j=1}^m\left \{\sum\limits_{i=1}^v A^i_{j-n}\right\}\right ]}{S_{t-


6}^{t-1}}\,,


    \label{e3-2-14}


    \end{equation}


где  $A^i_{j-n}$~--- ПАС $j$-го сотрудника в $i$-м журнале за $n$-й 


год, которая вы\-чис\-ля\-ет\-ся по формуле~(\ref{e2-2-14});


$v$~--- количество журналов в системе <<Web of Science>>;


$m$~--- численность сотрудников в соответствующей научной 


организации;


$t$~--- отчетный год, на который проводится оценка 


результативности деятельности научной организации;


$S_{t-6}^{t-1}$~--- общее число публикаций сотрудников научной 


организации в системе <<Web of Science>> за последние 


пять лет, начиная с года, предшествующего отчетному 


моменту времени.


    


    В формализованной канонической форме ИФПР выглядит еще более 


серь\-ез\-ным и значимым по сравнению с его текстовым описанием. Однако и 


каноническая форма, и текстовое описание~--- это все же туман, за которым 


скрывается сущность ИФПР с учетом конкретных обстоятельств. 


Принципиальное практическое обстоятельство, которое <<скрашивает>> 


витиеватые формулировки, состоит в количестве публикаций одного автора в 


том или ином журнале за два смежных года. Теоретически этот показатель 


может принимать произвольные (целые и положительные) значения. Однако 


выборочная проверка в нескольких журналах в области информатики, 


проведенная собственными силами, показала, что за период 2007--2009~гг.\ 


не нашлось ни одного автора, разместившего более одной статьи в одном и 


том же журнале системы <<Web of Science>> за два смежных года.


{ %\looseness=1





}


    


    Таким образом, если предположить, что каждый автор размещает в 


одном журнале не более одной статьи, то после замены $P^i_{j-n}$ 


в~(\ref{e2-2-14}) на <<единицу>> ПАС \mbox{$j$-го} сотрудника в $i$-м журнале за 


$n$-й год превращается в банальный им\-пакт-фак\-тор $i$-го журнала за 


$n$-й год и не более того. Произведя необходимую подстановку в~(\ref{e3-2-14}) 


и исключая суммирование по сотрудникам, получаем реальное 


значение ИФПР почти в 100\% случаев:





\noindent


    \begin{equation}


    \mathrm{ИФПР} = \fr{\sum\limits_{n=t-6}^{t-1}\left( \sum\limits_{i=1}^v 


k_n^i I_n^i\right)}{S_{t-6}^{t-1}}\,,


    \label{e4-2-14}


    \end{equation}


где $k_n^i$~---  количество публикаций, размещенных в $i$-м журнале 


за $n$-й год всеми сотрудниками научной организации;


$v$~--- количество журналов в системе <<Web of Science>>;


$t$~---   отчетный год, на который проводится оценка 


результативности деятельности научной организации;


$S_{t-6}^{t-1}$~---  общее число публикаций сотрудников научной 


организации в системе <<Web of Science>> за последние 


пять лет, начиная с года, предшествующего отчетному 


моменту времени.


    


    Что же означает такой многообещающий, на первый взгляд, показатель? 


Ровным счетом ничего, по крайней мере, что касается сформулированного в 


ТМО ИФПР. 


В~результате получается среднее значение традиционного им\-пакт-фак\-то\-ра 


нескольких журналов, т.\,е.\ сумма им\-пакт-фак\-то\-ров нескольких журналов, 


поделенная на количество статей, размещенных в этих журналах. Например, 


если в ка\-ком-то году размещено три статьи в трех журналах (по одной в 


каждом) с им\-пакт-фак\-то\-ра\-ми 0,5, 1,5 и 2,2, то ИФПР равен 1,4, как это 


представлено в табл.~1.


    


    Здесь <<всплывают>> и другие <<родовые пятна>> традиционного 


импакт-фак\-то\-ра, связанные с неоднородной популярностью разных журналов 


независимо от года наблюдения. Большинство научных организаций РАН 


узко профилированы, и трудно представить, что сотрудники института, 


за\-ни\-ма\-юще\-го\-ся проб\-ле\-ма\-ми ядерной энергетики, будут размещать свои 


публикации в журналах, посвященных маркшейдерскому делу или 


проблемам инерциальной навигации. Фактически это обстоятельство сводит 


ИФПР к среднему значению им\-пакт-фак\-то\-ра того или иного журнала за пять лет.





\begin{table}\small


\begin{center}


\Caption{Варианты расчета ИФПР}


\vspace*{2ex}





\begin{tabular}{|c|c|c|c|c|c|c|c|c|}


\hline


№ журнала&Импакт-фактор &\multicolumn{7}{c|}{Количество статей в каждом журнале}\\


\hline


1&0,5&1&2&0&\cellcolor[gray]{.8}0&2&3&\cellcolor[gray]{.8}5\\


2&1,5&1&2&1&\cellcolor[gray]{.8}0&1&1&\cellcolor[gray]{.8}2\\


3&2,2&1&2&1&\cellcolor[gray]{.8}1&1&1&\cellcolor[gray]{.8}1\\


\hline


\multicolumn{2}{|r|}{Всего статей }&3&6&2&\cellcolor[gray]{.8}1&4&5&\cellcolor[gray]{.8}8\\


\hline


\multicolumn{2}{|r|}{ИФПР}&1,4&1,4&1,85&\cellcolor[gray]{.8}2,2&1,175&1,04&\cellcolor[gray]{.8}0,96\\


\hline


\end{tabular}


\end{center}


\vspace*{6pt}


\end{table}





    Из расчетов, приведенных в табл.~1, видна абсурдность нового 


показателя. В~столбцах, выделенных заливкой, видно, что наибольшее 


значение ИФПР до-\linebreak стигается при минимальном количестве статей (одной), а 


наименьшее~--- при наи-\linebreak большем количестве статей (восьми). Конечно, 


<<мал золотник, да дорог>>, но ведь импакт-фактор, как известно, никак не 


определяет ценность статьи и изложенного в ней исследования.


{ %\looseness=-2





}


    


    Трудно понять, что имели в виду разработчики этого показателя и как 


они хотят его применять к категорированию научных организаций, особенно 


при отнесении их к третьей категории, предусматривающей такие серьезные 


меры, как ликвидация или смена руководителя.





\section{Заключение}


    


    Ошибки в учебниках опасны тем, что ученик воспринимает все 


написанное в учебнике как абсолютную истину и успевает наделать много 


ошибок, пока реалии практической деятельности не откроют ему глаза на 


недоработки, заложенные в учебник его авторами. Поэтому приходится 


соглашаться с необходимостью длительных проверок и экспертиз перед тем 


как учебник включат в учебный процесс. В~противном случае велика 


опасность того, что в учебный материал попадут ошибки. Правда, даже в 


этом случае сохраняется возможность исправить возникшие ошибки, если на 


этапе обучения их заметит квалифицированный преподаватель.


    


    Ошибки в нормативно-законодательном регулировании нашей жизни 


еще страшнее, так как у исполнителей остается только два выбора: один~--- это 


повторять навязанные ошибки, а второй~--- нарушать законы, нормы и 


правила. Что касается ИФПР, то наиболее целесообразно исключить его из 


перечня показателей для оценки результативности деятельности научных 


организации.





{\small\frenchspacing


{%\baselineskip=10.8pt


\addcontentsline{toc}{section}{Литература}


\begin{thebibliography}{9}








\bibitem{1-2-14}


Типовая методика оценки результативности деятельности научных 


организаций, выполняющих научно-исследовательские, 


опытно-конструкторские и технологические работы гражданского 


назначения (ТМО), введенная в действие приказом Министра 


образования и науки Российской Федерации от 14.10.2009 №\,406.





\bibitem{2-2-14}


Science Watch. {\sf http://sciencewatch.com/dr/sci/10/aug15-10\_1/}.





       \label{end\stat}





\bibitem{3-2-14}


Импакт-факторы некоторых российских журналов в 2009~году. {\sf 


http://www. physchem.chimfak.rsu.ru/General/if-2009rus.html}.


\end{thebibliography}


}


}         


     





     


